\section{Limitations and practical feasibility}
\label{sec:limitations}

The claims supported by this study are bounded by its physical model,
observation model, and finite registered design. We state those boundaries
explicitly before discussing field deployment.

\subsection{Physical model and label semantics}

\begin{itemize}
  \item \textbf{Simulation only.} All responses are generated by the registered
  TTBI model; no field measurement enters the study. Every conclusion is
  conditional on this simulated distribution.

  \item \textbf{Two-dimensional vertical dynamics.} The model captures
  vertical and pitch response only. Lateral, torsional, soil-continuum, and
  pier-tilt behaviour --- all potentially informative for scour --- are absent.

  \item \textbf{The support spring is not soil--structure interaction.} One
  vertical spring per bridge support represents the foundation. The target is
  modeled support-stiffness loss in percentage points, not scour depth,
  embedment loss, or eroded volume; no universal hydraulic-to-stiffness mapping
  is claimed.

  \item \textbf{Bearing and crack labels are design coordinates.} Nominal
  abutment rotational fixity is not a calibrated physical bearing-condition
  percentage. The crack arm is a uniform local finite-element $EI$ reduction,
  not a breathing crack or calibrated crack-depth model
  \cite{sinha2002simplified}.

  \item \textbf{Fixed profile and restricted nuisance scope.} One generated
  FRA-class-4 profile realization is shared by every state and passage. Rail
  profile is therefore not an operational EOV in Paper~1, and no claim of
  transport to another track section is available. Track-layer damage, wheel
  polygonization, and suspension damage are disabled. Evaluating the frozen
  models later under varied profile phases would be a distribution-shift
  experiment, not robustness to an EOV represented during training.

  \item \textbf{The nominal track is an inherited simplified topology.} The
  planar source combines two-rail quantities for rail inertia/mass and a full
  sleeper, but retains other values that Zhai et al.~\cite{zhai2004track}
  report per rail seat. It also retains source values derived at 0.545~m on a
  0.600~m support lattice, omits the adjacent-ballast $K_w,C_w$ shear branch,
  and condenses each on-bridge ballast lump onto the support-aligned deck node.
  The production baseline is therefore neither a demonstrated one-seat model
  nor a consistently doubled two-rail model, and it is not a complete Zhai
  reproduction. Prospective one-seat/two-rail and 0.545-to-0.600~m
  parameter-transfer sensitivities quantify two parameter-level choices; the
  omitted-shear/condensation topology itself remains unexamined. If a wheelset
  proxy wins the channel selection, that topology limitation must be resolved
  before publication or the wheelset result cannot be headlined.

  \item \textbf{Damping is model-form dependent.} The 0.1\% rail Rayleigh
  target is an inherited author-chosen value, not a property reported by Zhai;
  a prospective damping sensitivity bounds that choice. Production recomputes
  bridge coefficients for each realized state and recomputes them separately
  on every mesh in refinement studies, with rail damping referenced to the
  corresponding bridge modes. A fixed-healthy/fixed-grid coefficient arm is a
  model-form sensitivity. Consequently, stiffness interventions in the
  production model alter both assembled stiffness and Rayleigh-dependent
  damping; they are not evidence for physical damage-dependent damping.

  \item \textbf{Bilateral contact.} The linear wheel--rail model does not
  simulate separation and re-contact. The authenticated contact gate aborts
  generation on registered violations, but passing that gate does not quantify
  the response error caused by the bilateral approximation.

  \item \textbf{Author-chosen operational design.} The speed and temperature
  envelopes, deck-modulus temperature law, and vehicle-property perturbations
  were not fitted to a route, climate, or fleet population. Results are
  conditional on those finite design choices.
\end{itemize}

\subsection{Observation model and channel-level claims}

\begin{itemize}
  \item \textbf{Noise is a symmetric stress, not a hardware model.} The 5\%
  relative-noise arm applies the same mathematical rule to every channel and
  includes no transducer noise density, bias stability, bandwidth, mounting
  transfer function, or environmental qualification.

  \item \textbf{Channels are not sensors.} The experiment ranks modeled
  responses, not physical transducer counts or placements. In particular, the
  two \texttt{physical8\_v1} wheelset rows are total constrained-motion
  acceleration proxies assembled from the linear model. They omit axle-box
  mounting dynamics, local wheel/rail compliance beyond the model, sensor
  bandwidth and filtering, and contact nonlinearity. They must not be described
  as measured axle-box accelerations without a separate observation model.
\end{itemize}

\subsection{Statistical scope}

\begin{itemize}
  \item \textbf{Conditional finite-budget selection.} ``Best'' means best
  under the registered 16-cell search space, HPO budgets, fixed folds, seed
  lists, and F40-S channel screen. Equal budgets do not eliminate finite-search
  optimization error, and no global architecture or sensor optimum is claimed.

  \item \textbf{Development adjudication is selection, not stability.} The
  480 Option-C out-of-fold fits choose the F40-S vectors and therefore cannot
  estimate post-selection stability. Stability is reported separately from 30
  disjoint post-freeze initialization seeds per retained pipeline and block on
  the sealed outer test; those runs cannot re-rank a model.

  \item \textbf{The channel screen has an asymmetric role.} Eight singles are
  reported but non-selectable; only the 28 pairs can win. Subsequent HPO on the
  selected pair may improve its block-local model but is not evidence that the
  pair beats the 27 pairs that were screened under frozen F40-S settings.

  \item \textbf{Descriptive uncertainty, not population coverage.} State-first
  resampling and across-seed summaries are finite-design sensitivity analyses.
  They are not confidence intervals for a field population, familywise-error
  controlled tests, or significance/superiority guarantees.

  \item \textbf{Independent tuning changes the cross-block estimand.} Primary
  results compare pipelines after block-local HPO. They do not estimate
  zero-shot transfer. Frozen F40-S settings are evaluated downstream only in a
  separately labelled, secondary non-selection analysis.

  \item \textbf{No detection claim.} The study reports continuous regression
  error and localization. Sensitivity, specificity, probability of detection,
  calibrated damage probability, and minimum detectable severity would require
  a development-selected decision threshold that this protocol does not define.
\end{itemize}

\subsection{Practical feasibility and implementation challenges}

Field use first requires a hardware-specific observation model. Transducers,
mounts, acquisition chains, and environmental loads differ markedly across a
vehicle; axle-box environments are especially demanding~\cite{en61373}.
Second, supervised deployment requires trustworthy labels or a credible route
through model updating, domain adaptation, or hybrid physics--data methods.
The simulation-to-field gap --- model-form error, unmodeled EOVs, traffic
variability, and asset variability --- remains the dominant unquantified risk
\cite{malekjafarian2015review, fernandes2026canelas}. Third, the spatial
transform assumes known passage speed and repeated passages require fleet and
positioning integration. Finally, the models cover one train class and two
idealized bridge geometries. Transfer across vehicles, bridge types, profiles,
and speeds outside 70--90~km/h is untested. A staged pathway should therefore
begin with instrumented healthy crossings for EOV and observation-model
calibration, followed by controlled stiffness interventions or monitored
scour-critical structures.
